\clearpage
\section{Electrical Model and Power Accounting}
\label{sec:electrical}
\subsection{Conventions and state variables}
A reproducible plant model requires consistent device bases and connection signs.
Each supply block consists of a unit transformer, an active front end (AFE), a voltage source inverter (VSI), a power supply unit (PSU) array and a server DC--DC stage.
The unit transformer belongs to the algebraic network, and the four converter stages carry the differential states of the block.
Each supply block uses its own three-phase rating $S_n$ as the power base.
The network uses $S_b=200$~MVA and $\omega_b=2\pi60$~rad/s.
Time is in seconds.
Per-unit inductances and capacitances multiply derivatives divided by $\omega_b$.
Controller integrators retain the unscaled time derivatives used in the source equations.
The reference speed is $\omega_s=1$.

The local AFE frame rotates by the phase-locked loop (PLL) angle $\theta$.
Define
\begin{equation}
\R(\theta)=\begin{bmatrix}\cos\theta&\sin\theta\\-\sin\theta&\cos\theta\end{bmatrix},\qquad
\J=\begin{bmatrix}0&-1\\1&0\end{bmatrix}.
\end{equation}
Here $\bm v_g$ is the voltage at the low-voltage terminal of the unit transformer in network coordinates.
The rectifier current $\bm i_a$ is positive into the supply block.
The VSI uses a separate synchronous frame with voltage $\bm v_v$ and converter current $\bm i_c$.

One supply block has 21 differential states.
The AFE contributes $\theta$, $\epsilon$, $\bar v_q$, the two entries of $\bm i_a$, $\xi_a$ and the two entries of $\bm\gamma_a$.
The VSI contributes the two entries of each of $\bm i_c$, $\bm v_v$, $\bm\xi_v$ and $\bm\gamma_v$.
The remaining states are the shared DC-link voltage $v_d$, the PSU output voltage $v_p$ and integrator $\xi_p$, and the server voltage $v_e$ and integrator $\xi_e$.
The following vector equations count each two-entry vector as two states.
The accompanying paper writes $v_{\rm dc}$, $v_{\rm psu}$ and $v_{\rm s}$ for $v_d$, $v_p$ and $v_e$.

\subsection{AFE rectifier}
The rectifier synchronizes to the unit-transformer voltage and regulates the shared DC link.
Let $\bm v_a=\R(\theta)\bm v_g$.
Its equations are
\begin{align}
\dot\theta&=\omega_b(\omega_{\rm pll}-\omega_s), & \dot\epsilon&=\bar v_q,\\
\dot{\bar v}_q&=\omega_{\rm lp}(v_{a,q}-\bar v_q), &
\omega_{\rm pll}&=\omega_s+k_{p,\rm pll}\bar v_q+k_{i,\rm pll}\epsilon,\\
\frac{\ell_a}{\omega_b}\dot{\bm i}_a&=\bm v_a-v_d\bm m_a-r_a\bm i_a+\omega_{\rm pll}\ell_a\J\bm i_a,\\
\dot\xi_a&=v_d^\star-v_d, & \dot{\bm\gamma}_a&=\bm i_a-\bm i_a^\star,\\
\bm i_a^\star&=\begin{bmatrix}k_{p,da}(v_d^\star-v_d)+k_{i,da}\xi_a\\0\end{bmatrix},\\
\bm m_a&=\frac{k_{p,ca}(\bm i_a-\bm i_a^\star)+k_{i,ca}\bm\gamma_a+\omega_{\rm pll}\ell_a\J\bm i_a}{v_d}.
\end{align}
A superscript $\star$ denotes a reference value.
The constants $r_a$ and $\ell_a$ are the AFE filter resistance and inductance.
Subscripts $da$ and $ca$ identify the DC-voltage and current controllers.
The current-controller error and the coupling signs follow the source model \cite{zhao2026dynamic}.
Another rectifier convention would change the model.

\subsection{VSI and shared DC link}
The VSI regulates its output voltage while drawing energy from the rectifier DC link.
With reference $\bm v_v^\star=[v_u^\star,0]^{\mathsf T}$, its equations are
\begin{align}
\frac{\ell_v}{\omega_b}\dot{\bm i}_c&=v_d\bm m_v-\bm v_v-r_v\bm i_c+\omega_s\ell_v\J\bm i_c,\label{eq:vsi-current}\\
\frac{c_v}{\omega_b}\dot{\bm v}_v&=\bm i_c-\bm i_o+\omega_s c_v\J\bm v_v,\\
\dot{\bm\xi}_v&=\bm v_v^\star-\bm v_v, & \dot{\bm\gamma}_v&=\bm i_c^\star-\bm i_c,\\
\bm i_c^\star&=k_{p,vv}(\bm v_v^\star-\bm v_v)+k_{i,vv}\bm\xi_v-\omega_s c_v\J\bm v_v,\\
\bm m_v&=\frac{k_{p,cv}(\bm i_c^\star-\bm i_c)+k_{i,cv}\bm\gamma_v-\omega_s\ell_v\J\bm i_c}{v_d},\label{eq:vsi-modulation}\\
\frac{c_d}{\omega_b}\dot v_d&=\bm m_a^{\mathsf T}\bm i_a-\bm m_v^{\mathsf T}\bm i_c.
\end{align}
Here $r_v$, $\ell_v$ and $c_v$ are the VSI filter parameters, $c_d$ is the DC-link capacitance, and $\bm i_o$ is the current drawn by the PSU array.
Subscripts $vv$ and $cv$ denote its voltage and current controllers.
The DC-link equation couples both converters through their current balance.
The study represents online double-conversion operation.

\subsection{PSU array and server DC--DC stage}
The positive-sequence PSU equivalent converts its AC input into a common DC supply.
Let $g_p$ be its controlled input conductance and $i_p$ its DC output current.
Then
\begin{align}
\frac{c_p}{\omega_b}\dot v_p&=\frac{(g_p-r_p g_p^2)\bm v_v^{\mathsf T}\bm v_v}{3v_p}-i_p, &
\dot\xi_p&=v_p^\star-v_p,\\
g_p&=k_{p,vp}(v_p^\star-v_p)+k_{i,vp}\xi_p, & \bm i_o&=g_p\bm v_v.
\end{align}
The parameters $r_p$ and $c_p$ represent the PSU equivalent resistance and capacitance.
The downstream DC--DC stage has a fast current loop represented by its commanded current $i_e$, with
\begin{align}
\frac{c_e}{\omega_b}\dot v_e&=i_e-g_e v_e, & \dot\xi_e&=v_e^\star-v_e,\\
g_e&=\frac{p_{\rm load}}{3(v_e^\star)^2}, &
i_e&=k_{p,ve}(v_e^\star-v_e)+k_{i,ve}\xi_e, &
i_p&=\frac{v_e}{v_p}i_e.
\end{align}
The factor 3 follows the source model's three-phase power convention.
The input $p_{\rm load}$ is the prescribed server power on the block base at $v_e=v_e^\star$.
The instantaneous power of the modeled load is $3g_e v_e^2$, so it equals $p_{\rm load}$ at the voltage reference.
This distinction matters when interpreting internal voltage excursions.
The server power of the hall sets this command.
It is not an additional current injection at the PCC.

\subsection{Hall and network connections}
Each hall has one common feeder and four parallel supply blocks.
Each block has its own 3~MVA unit transformer and converters.
The four blocks of hall $h$ have identical parameters, and each receives the command $p_{\rm load}=P_{{\rm IT},h}/(4S_n)$, where $P_{{\rm IT},h}$ is the server power of the hall defined in Section~\ref{sec:allocation}.
The AC input power of a block is the active power that it draws at the low-voltage terminal of its unit transformer, and $P_{{\rm AC},h}$ denotes its sum over the four blocks of hall $h$.
The cooling load connects to the hall bus after the common feeder and before the unit transformers.
Auxiliary loads connect to the two medium-voltage buses.
The full model keeps both medium-voltage buses, both main transformers, the 12 feeders and the 48 unit transformers as separate elements.
Section~\ref{sec:equivalent} describes how the single-block equivalent merges them.

The current drawn from the network by supply block $n$ is
\begin{equation}
\bm i_{n,\rm sys}=\frac{S_n}{S_b}\R(\theta_n)^{\mathsf T}\bm i_{a,n}.
\end{equation}
With the network voltage phasors $V$, generator injections $I_G$ and explicit load withdrawals $I_L$, nodal balance is $YV=I_G-I_L$.
A load represented as an admittance is included in $Y$ and is excluded from $I_L$.
The network is algebraic and retains the line and transformer impedances and the feeder charging.

In the full model, the cooling power of a hall at the operating point is 17\% of the AC input power of its supply blocks.
The total auxiliary power is 3\% of the AC input power of all supply blocks, divided equally between the medium-voltage sections.
Cooling uses a constant-power representation and the auxiliaries use constant impedance.
With the AC input power $P_{n,\rm AC}$ of block $n$ at the operating point, the static demand at bus~8 is
\begin{equation}
P_{8,\rm residual}=100-1.20\sum_n P_{n,\rm AC}\quad\mathrm{MW}.
\end{equation}
The original 35~Mvar remains as a static impedance load.
Collector losses are accounted for by the network.
The PCC power is the sum of the powers entering the main transformers, so it includes these losses.
The single-block equivalent keeps the static loads of the full model.
They are not adjusted to cancel the loss error of the equivalent.

The external system has a fourth-order synchronous machine with excitation and turbine-governor dynamics at bus~1, a droop grid-forming converter at bus~2 and a phase-locked-loop grid-following converter at bus~3.
Their equations and parameter conventions follow the modified nine-bus reference in the supplement of \cite{zhao2026dynamic}.
The implementation uses the source-current feed-forward sign $-1$ and the q-axis voltage convention for the grid-following converter.
The accompanying data files give all grid parameter values and network branches.
These grid devices contribute 37 states, so a model with $N$ supply blocks has $37+21N$ states.
The full model has $N=48$ blocks and 1045 states, and the single-block equivalent has $N=1$ and 58 states.
